\section{CONSIDERAÇÕES FINAIS}
\label{sec:consideracoes_finais}

A equipe percorreu duas etapas do desenvolvimento. Na primeira, de definição do produto, tratou da definição do mercado, da identificação das necessidades dos usuários, do desdobramento dessas necessidades em requisitos técnicos e da elaboração dos desenhos iniciais. Na segunda, de projeto conceitual, tratou da modelagem funcional, do estudo de diferenciação, da escala vertical e do valor mercadológico, do estudo de aproveitamento técnico, da reformulação dos desenhos e do delineamento da comercialização.

A definição do mercado delimitou como segmento-alvo as redes de moda \textit{mass market} e \textit{fast fashion} que já utilizam self-checkout, núcleo formado por Renner, Riachuelo e C\&A. Dentro desse segmento, a expansão de 15 a 20 lojas conceito prevista pela Riachuelo para 2026 foi identificada como o mercado obtível no curto prazo, por coincidir com o momento em que a rede padroniza sua tecnologia de checkout.

A pesquisa com consumidores, realizada com um survey de 142 respostas e onze entrevistas, mostrou que o comportamento de compra se organizou menos pelo perfil sociodemográfico ou pelo tipo de loja frequentada e mais por duas dimensões: a orientação para autonomia, tecnologia e eficiência e a dor de disponibilidade de produto. A partir delas, a clusterização identificou quatro grupos, com reações muito diferentes ao conceito de compra sem checagem manual: aceitação total entre os autônomos digitais, rejeição quase unânime entre os dependentes de atendimento e neutralidade predominante nos demais. No conjunto da amostra, apenas 24\% endossaram esse processo. As entrevistas acrescentaram que a preferência por autonomia varia conforme a etapa da jornada e que, para quem depende de ajuda, perder o contato com um vendedor pesa mais do que o risco de falha técnica.

Esses achados deram origem a oito necessidades, sete desejos e cinco demandas. Sete dessas vozes foram desdobradas em requisitos técnicos por meio do QFD, e a liberação da peça condicionada à confirmação do pagamento, o tempo de pagamento e a clareza do fluxo no aplicativo concentraram cerca de 71\% da importância técnica. Na matriz de correlação, o acesso a atendimento humano apareceu em conflito com os requisitos de rapidez e autonomia, conflito que a solução precisará acomodar. Em seguida, a equipe traduziu esses requisitos no esquema à mão livre de um dispositivo rígido de 98 por 82 por 26 mm, com retenção por embreagem de esferas e liberação autorizada pelo celular do cliente, sem depender da rede da loja no momento da saída.

No projeto conceitual, a modelagem funcional definiu a função total do sistema, conciliar peça e pagamento, e a desdobrou em três funções básicas, treze secundárias necessárias e cinco acessórias. O estudo de diferenciação comparou seis soluções existentes e mostrou que a função de liberar uma trava física pelo celular já foi oferecida no mercado. A diferenciação do produto ficou, por isso, na combinação entre retenção física, compra no salão, uso da identificação por radiofrequência que a Riachuelo já tem e atendimento humano em cada falha. A escala vertical, aplicada a três profissionais do varejo, indicou um valor mercadológico entre R\$~5 e R\$~8 por etiqueta para redes de médio e grande porte, cerca de uma ordem de grandeza abaixo do custo estimado do protótipo, o que reforçou o fornecimento sob locação.

O estudo de aproveitamento técnico reuniu, em uma galeria de ideias, 101 efeitos físicos e 126 portadores para as 18 funções, com referências de seis linhas de similaridade. Na reformulação dos desenhos, a matriz morfológica combinou os princípios de sete funções em três alternativas de concepção, e a matriz de decisão apontou a segunda. Ela mantém o pino com embreagem de esferas do desenho inicial e substitui o atuador por uma mola armada, solta por um gatilho. A autorização passa a ocorrer pelo toque do celular, a liberação é indicada por um anel de LEDs e a bateria é recarregada por indução. A comercialização foi delineada como venda direta à rede, com contrato de serviço, e a primeira implantação foi proposta para uma única loja.

Em conjunto, a equipe concluiu que a solução tende a gerar valor para um público definido, o dos clientes que já compram de forma autônoma e se incomodam com fila, e que seu alcance depende de combinar a liberação autônoma com a preservação do atendimento humano. Para os demais grupos, os encaixes identificados permanecem como hipóteses: a redistribuição do tempo dos vendedores, para os dependentes de atendimento; a informação de disponibilidade, para os frustrados com disponibilidade; e uma via rápida opcional, para os pragmáticos equilibrados.

Esses resultados têm limites. A amostra é pequena, e a forma de recrutamento dos participantes não foi documentada em detalhe, de modo que as proporções descrevem esta amostra, e não o conjunto de clientes da Riachuelo. Além disso, a operação e a gestão da loja, que também usarão a solução ou decidirão sobre ela, ainda não foram investigadas. No projeto conceitual, a clínica de preços teve três respostas, contra as 15 a 30 que o método prevê, as notas da matriz de decisão são avaliações da equipe, sem ensaio, e o desenho reformulado não foi dimensionado.

No projeto detalhado, a equipe deverá modelar o dispositivo em CAD, com cotas, materiais e a lista de materiais inicial, verificar se o mecanismo de mola e gatilho cabe no envelope especificado e confirmar com a Riachuelo o protocolo da etiqueta de radiofrequência e a tecnologia do antifurto da saída. Deverá também testar com os próprios clientes os encaixes entre a solução e cada grupo, verificar nos testes do protótipo as especificações-meta definidas para cada requisito e consolidar as dimensões do mecanismo de garra e a escolha do atuador a partir de medições em etiquetas comerciais. A substituição do carretel por uma peça não ferromagnética, que tornaria o dispositivo imune ao destravador magnético comercial, também permanece em avaliação.
